% ============================================
% LUC-Escolha (EN)
% English translation — SPEC-935-R355
% ============================================
\section*{\textbf{LUC-Escolha} --- The Choice}
\markboth{LUC-Escolha --- The Choice}{}

\noindent \textit{What to do with the devil you carry}

\noindent \textbf{April 2026 --- Barbacena}

\vspace{0.5em}\noindent\rule{\textwidth}{0.5pt}\vspace{0.5em}

\noindent Lúcia spent the night in Barbacena's only hotel. A small room, mold on the walls, a ceiling fan that spun like a slow carousel.

\noindent She placed the two diaries on the table: the original (brown cover, 62 years of writing) and the second one (black cover, written after death).

\noindent She did not open the second one. Not yet.

\noindent Instead, she sat on the bed and thought.

\noindent She had three options.

\noindent \textbf{Option 1: Return the diary to the archive.}

\noindent Pretend nothing had happened. Let another researcher find it. The cycle would continue --- but she would not be responsible. Or would she? She had already read it. The contamination had already begun. Returning the diary would not erase the symptoms.

\noindent \textbf{Option 2: Destroy the diary.}

\noindent Burn it. Tear it. Dissolve it in acid. If the physical medium were destroyed, the cycle might stop. But what about the second diary? And what if destroying the original diary freed the creature instead of imprisoning it?

\noindent \textbf{Option 3: Accept.}

\noindent Accept that the cycle was now part of her. Accept that the right eye was changing. Accept that the smell would not go away. Accept that she was, from that moment on, the diary's new keeper --- and that, sooner or later, she would have to pass it on.

\noindent \textbf{--- But there was a fourth option.}

\noindent Read the second diary. Not as a reaction to fear --- as an act of investigation. Actively choose to enter the cycle instead of being dragged into it. Refuse the role of victim and take on the role of witness. The difference, Lúcia thought, was the difference between drowning and diving.

\vspace{0.5em}\noindent\rule{\textwidth}{0.5pt}\vspace{0.5em}

\noindent The answer came at 3:14 in the morning.

\noindent The smell woke her.

\noindent It was not the sweetish smell from before. It was stronger. Denser. Smell of wet earth. Of burnt bone. Of something coming out of the ground.

\noindent Lúcia opened her eyes.

\noindent The second diary was open on the table.

\noindent She did not remember opening it.

\noindent The page showed a sentence, written in the same trembling hand:

\noindent \textit{\textquotedblleft{}She has already chosen for you. You just do not know it yet.\textquotedblright{}}

\noindent Lúcia closed the diary. Put both in her backpack. Paid the hotel. Took the 5 AM bus back to Belo Horizonte.

\noindent During the whole trip, her right eye burned.

\noindent As if something were growing inside it.

\noindent As if something were \textit{learning to see}.

\vspace{0.5em}\noindent\rule{\textwidth}{0.5pt}\vspace{0.5em}

\vspace{0.5em}
\noindent\textit{\textbf{↪ Routes:} \rota{LUC-06} \rotasep \rota{LUC-Entrevista}}

\vspace{0.5em}\noindent\rule{\textwidth}{0.5pt}\vspace{0.5em}

\subsection*{CONTAMINATION SCALE --- UPDATE 1}

\noindent \textit{The archive has detected that you followed Lúcia's investigation up to the moment of the choice.}

\begin{moltabletwo}
\hline
Points & Diagnosis \\
\hline
0--1 & Reader intact. Unlikely, but possible. \\
\hline
2--4 & Reader exposed. You have already read documents you should not have. \\
\hline
5--7 & Reader contaminated. Lúcia chose. You have made your choice too --- even if you do not know which one it was. \\
\hline
8--10 & Patient 1,263 forming. The creature already has your name. \\
\hline
\end{moltabletwo}

\begin{quote}
Archive note: the points are not cumulative. The scale adjusts as the archive closes around you.
\end{quote}
